% !TeX program = pdflatex
%==============================================================================
%  Lista Teórica E2 --- Redução de Dimensionalidade (PCA e t-SNE)
%  O gabarito sai de "Lista teorica E2 - gabarito.tex".
%==============================================================================
\providecommand{\gabaritoopt}{}
\documentclass[11pt,a4paper]{article}
\usepackage[\gabaritoopt]{../../recursos/latex/estilo-lista}

\begin{document}
\cabecalholista{E2}{Redução de Dimensionalidade (PCA e t-SNE)}

%------------------------------------------------------------------------------
\begin{exercicio}[PCA na mão]
Quatro pontos em $\R^2$:
\[
  (3,1), \quad (1,3), \quad (-3,-1), \quad (-1,-3).
\]
\begin{enumerate}[label=(\alph*),leftmargin=1.1cm]
  \item Verifique que a média é a origem e calcule a matriz de covariância
        amostral $\bm S = \tfrac1n\bm X^\top\bm X$.
  \item Encontre os autovalores e autovetores de $\bm S$. Escreva as duas
        componentes principais como vetores unitários.
  \item Qual a proporção da variância explicada pela primeira componente?
  \item Calcule as coordenadas dos quatro pontos no novo sistema (isto é,
        projete cada um sobre as duas componentes). Que estrutura geométrica isso
        revela?
\end{enumerate}
\end{exercicio}

\begin{solucao}
\textbf{(a)} A soma das coordenadas é
$(3+1-3-1,\ 1+3-1-3) = (0,0)$, logo a média é a origem e os dados já estão
centrados. Então
\[
  \bm X^\top\bm X =
  \begin{pmatrix} 9+1+9+1 & 3+3+3+3 \\ 3+3+3+3 & 1+9+1+9 \end{pmatrix}
  = \begin{pmatrix} 20 & 12 \\ 12 & 20 \end{pmatrix},
  \qquad
  \bm S = \frac14 \bm X^\top\bm X = \begin{pmatrix} 5 & 3 \\ 3 & 5 \end{pmatrix}.
\]

\smallskip
\textbf{(b)} Para uma matriz da forma $\begin{psmallmatrix}a&b\\b&a\end{psmallmatrix}$
os autovetores são sempre $(1,1)$ e $(1,-1)$, com autovalores $a+b$ e $a-b$.
Conferindo:
\[
  \bm S\begin{pmatrix}1\\1\end{pmatrix} = \begin{pmatrix}8\\8\end{pmatrix} = 8\begin{pmatrix}1\\1\end{pmatrix},
  \qquad
  \bm S\begin{pmatrix}1\\-1\end{pmatrix} = \begin{pmatrix}2\\-2\end{pmatrix} = 2\begin{pmatrix}1\\-1\end{pmatrix}.
\]
Logo $\lambda_1 = 8$ e $\lambda_2 = 2$, com componentes principais
\[
  \bm v_1 = \frac{1}{\sqrt2}(1,1) \approx (0{,}7071;\ 0{,}7071),
  \qquad
  \bm v_2 = \frac{1}{\sqrt2}(1,-1) \approx (0{,}7071;\ -0{,}7071).
\]
(Sinal e ordem: os autovetores são definidos a menos de sinal, e a convenção é
ordená-los por autovalor decrescente.)

\smallskip
\textbf{(c)} $\dfrac{\lambda_1}{\lambda_1+\lambda_2} = \dfrac{8}{10} = \mathbf{80\%}$.

\smallskip
\textbf{(d)} Projetando, $z_j = \x^\top\bm v_j$:
\begin{center}\small
\renewcommand{\arraystretch}{1.25}
\begin{tabular}{@{}lcc@{}}
\toprule
ponto & $z_1$ (sobre $\bm v_1$) & $z_2$ (sobre $\bm v_2$) \\
\midrule
$(3,1)$   & $4/\sqrt2 = 2\sqrt2 \approx 2{,}8284$ & $2/\sqrt2 = \sqrt2 \approx 1{,}4142$ \\
$(1,3)$   & $2\sqrt2$   & $-\sqrt2$ \\
$(-3,-1)$ & $-2\sqrt2$  & $-\sqrt2$ \\
$(-1,-3)$ & $-2\sqrt2$  & $\sqrt2$ \\
\bottomrule
\end{tabular}
\end{center}
A estrutura fica evidente: os quatro pontos formam um \textbf{retângulo}
centrado na origem, com o lado maior ($2\times 2\sqrt2 = 5{,}66$) na direção
$(1,1)$ e o menor ($2\times\sqrt2 = 2{,}83$) na direção $(1,-1)$. No sistema
original isso estava escondido, porque o retângulo está a $45^\circ$; a PCA
apenas girou os eixos para alinhá-los aos lados.

É a leitura geométrica da PCA em uma frase: \emph{ela procura o sistema de
coordenadas em que a nuvem tem eixos alinhados}, e a variância ao longo de cada
novo eixo é o autovalor correspondente. Aqui a razão entre os lados é
$2\sqrt2/\sqrt2 = 2$, e a razão entre os desvios-padrão é
$\sqrt{8}/\sqrt{2}=2$ --- consistente.
\end{solucao}

%------------------------------------------------------------------------------
\begin{exercicio}[quantas componentes guardar]
Uma matriz de dados com $d=8$ covariáveis tem autovalores da matriz de
covariância
\[
  \lambda = (12{,}5;\ 6{,}2;\ 3{,}1;\ 1{,}4;\ 0{,}9;\ 0{,}5;\ 0{,}3;\ 0{,}1).
\]
\begin{enumerate}[label=(\alph*),leftmargin=1.1cm]
  \item Calcule a proporção de variância explicada por cada componente e a
        proporção acumulada.
  \item Quantas componentes são necessárias para reter ao menos $90\%$ da
        variância? E $99\%$?
  \item O que significa \texttt{PCA(n\_components=0.90)} no
        \texttt{scikit-learn}, e por que é uma forma mais conveniente de
        especificar a redução do que dar o número de componentes?
  \item A soma dos autovalores vale $25$. A que quantidade dos dados originais
        esse número corresponde? O que isso diz sobre a necessidade de
        \textbf{padronizar} antes da PCA?
\end{enumerate}
\end{exercicio}

\begin{solucao}
\textbf{(a)} A soma é $\sum_j\lambda_j = 25$.
\begin{center}\small
\renewcommand{\arraystretch}{1.25}
\begin{tabular}{@{}lcccccccc@{}}
\toprule
componente & 1 & 2 & 3 & 4 & 5 & 6 & 7 & 8 \\
\midrule
$\lambda_j$        & $12{,}5$ & $6{,}2$ & $3{,}1$ & $1{,}4$ & $0{,}9$ & $0{,}5$ & $0{,}3$ & $0{,}1$ \\
proporção          & $0{,}500$ & $0{,}248$ & $0{,}124$ & $0{,}056$ & $0{,}036$ & $0{,}020$ & $0{,}012$ & $0{,}004$ \\
acumulada          & $0{,}500$ & $0{,}748$ & $0{,}872$ & $\mathbf{0{,}928}$ & $0{,}964$ & $0{,}984$ & $\mathbf{0{,}996}$ & $1{,}000$ \\
\bottomrule
\end{tabular}
\end{center}

\smallskip
\textbf{(b)} Para $90\%$: três componentes dão $87{,}2\%$, ainda insuficiente;
com quatro chega-se a $92{,}8\%$. São necessárias \textbf{4 componentes}. Para
$99\%$: sete componentes dão $99{,}6\%$ e seis dão $98{,}4\%$; são necessárias
\textbf{7}.

Vale reparar no rendimento decrescente: as $4$ primeiras componentes carregam
$92{,}8\%$, e as $4$ últimas somam $7{,}2\%$. Guardar metade das colunas custa
menos de $8\%$ da variância.

\smallskip
\textbf{(c)} Quando \texttt{n\_components} é um \emph{float} em $(0,1)$, o
\texttt{scikit-learn} o interpreta como a \textbf{proporção de variância a
reter}, e escolhe automaticamente o menor número de componentes que a atinge.
Neste exemplo, \texttt{PCA(n\_components=0.90)} guardaria $4$.

É mais conveniente porque $0{,}90$ é uma quantidade \textbf{interpretável e
transferível}: significa a mesma coisa em qualquer conjunto de dados, enquanto
``4 componentes'' depende de $d$ e da estrutura de correlação. Num
\texttt{Pipeline} com validação cruzada, o número escolhido pode até variar entre
dobras --- e é isso que se quer, já que o critério é o mesmo.

\smallskip
\textbf{(d)} A soma dos autovalores é o \textbf{traço} de $\bm S$, isto é, a soma
das variâncias das $d$ covariáveis originais --- a ``variância total'' dos dados.
(Girar os eixos não muda o traço.)

E aí está o problema: essa soma é medida nas \textbf{unidades} das covariáveis.
Se uma delas estiver em milímetros e outra em quilômetros, a primeira terá
variância $10^{12}$ vezes maior e sozinha responderá por praticamente toda a
variância total. A primeira componente principal será, essencialmente, essa
covariável --- não porque ela importa, mas porque a régua dela é menor.

Por isso, salvo quando todas as covariáveis já estão na mesma unidade e a escala
relativa entre elas é significativa, \textbf{padroniza-se antes da PCA} --- o que
equivale a fazer a decomposição da matriz de \emph{correlação} em vez da de
covariância. É o mesmo argumento do Exercício~2(c) da Lista Teórica~02, sobre
Ridge e Lasso.
\end{solucao}

%------------------------------------------------------------------------------
\begin{exercicio}[o que o t-SNE preserva, e o que não preserva]
O t-SNE constrói, no espaço original, probabilidades $p_{ij}$ de que $i$ escolha
$j$ como vizinho, e procura pontos $\bm z_i$ em $\R^2$ cujas probabilidades
$q_{ij}$ minimizem $\mathrm{KL}(p\,\|\,q)$.
\begin{enumerate}[label=(\alph*),leftmargin=1.1cm]
  \item A divergência de Kullback--Leibler $\sum_{ij} p_{ij}\log(p_{ij}/q_{ij})$
        não é simétrica. Que tipo de erro ela pune muito, e que tipo ela quase
        não pune? \emph{Sugestão:} olhe o que acontece com o somando quando
        $p_{ij}$ é grande e $q_{ij}$ pequeno, e vice-versa.
  \item Conclua o que se pode e o que não se pode ler num diagrama de t-SNE
        sobre: (i) pontos que aparecem próximos; (ii) a distância entre dois
        grupos distantes; (iii) o tamanho relativo de dois grupos.
  \item O parâmetro \texttt{perplexity} controla, grosso modo, quantos vizinhos
        cada ponto considera. O que se espera do diagrama com
        \texttt{perplexity} muito pequena? E muito grande?
  \item Por que não se usa t-SNE como etapa de um \texttt{Pipeline} de predição,
        do jeito que se usa PCA?
\end{enumerate}
\end{exercicio}

\begin{solucao}
\textbf{(a)} O somando é $p_{ij}\log(p_{ij}/q_{ij})$, e o peso de cada termo é
$p_{ij}$ --- a probabilidade no espaço \emph{original}.

Se $p_{ij}$ é grande (os pontos são vizinhos de verdade) e $q_{ij}$ é pequeno
(ficaram longe no mapa), o logaritmo vai a $+\infty$ e o termo é multiplicado por
um $p_{ij}$ grande: \textbf{punição enorme}. Separar vizinhos verdadeiros é
caríssimo.

Se $p_{ij}$ é pequeno (os pontos são distantes) e $q_{ij}$ é grande (ficaram
perto no mapa), o logaritmo vai a $-\infty$, mas o termo é multiplicado por
$p_{ij} \approx 0$: \textbf{punição quase nula}. Aproximar pontos distantes é de
graça.

Em uma frase: a KL nesta direção pune \emph{omitir} vizinhança verdadeira e é
quase indiferente a \emph{inventar} vizinhança falsa entre pontos distantes.

\smallskip
\textbf{(b)}
\begin{enumerate}[label=(\roman*),leftmargin=1.1cm]
  \item \textbf{Pontos próximos no mapa: pode confiar.} Como separar vizinhos
        verdadeiros é caríssimo, dois pontos que aparecem juntos quase sempre
        eram vizinhos no espaço original. É exatamente o que o t-SNE otimiza.
  \item \textbf{Distância entre grupos distantes: não significa nada.} Se dois
        grupos estão longe no espaço original, todos os $p_{ij}$ entre eles são
        praticamente zero, e a KL é indiferente a onde eles vão parar no mapa.
        Dizer ``o grupo A está mais perto do B do que do C'' a partir de um
        diagrama de t-SNE é uma leitura sem base.
  \item \textbf{Tamanho relativo dos grupos: também não.} O t-SNE normaliza a
        escala local de cada ponto pela \texttt{perplexity}, de modo que uma
        nuvem densa e uma esparsa podem ocupar áreas parecidas no mapa. A área de
        um agrupamento no diagrama não é proporcional a nada.
\end{enumerate}

\smallskip
\textbf{(c)} Com \texttt{perplexity} muito pequena, cada ponto só considera
pouquíssimos vizinhos: o mapa se fragmenta em muitos grupinhos minúsculos, e a
estrutura de médio alcance desaparece. Com \texttt{perplexity} muito grande, cada
ponto considera quase todo mundo, os $p_{ij}$ ficam quase uniformes e o mapa
tende a uma nuvem única e sem estrutura --- perto do que o PCA daria.

Vale reter que \texttt{perplexity} não é um parâmetro que se ``otimiza'': não há
função de risco a minimizar. Rodar com dois ou três valores e ver o que
\emph{sobrevive} aos três é a conduta usual.

\smallskip
\textbf{(d)} Por duas razões, e a primeira é decisiva.

\textbf{Não existe \texttt{transform}.} O t-SNE não aprende um mapa
$\R^d\to\R^2$; ele resolve um problema de otimização cujas variáveis são as
\emph{posições dos pontos que você lhe deu}. Não há como projetar uma observação
nova sem refazer o ajuste --- e refazê-lo com o ponto novo muda as posições de
todos os outros. Um `Pipeline` precisa de um passo que possa ser ajustado no
treino e aplicado à validação, e o t-SNE não oferece isso. (Na API do
\texttt{scikit-learn} isso aparece como um \texttt{fit\_transform} sem
\texttt{transform} correspondente.)

\textbf{E o objetivo é outro.} O t-SNE existe para produzir uma figura que um
humano olha. As duas coordenadas que ele devolve não são features com significado
--- não são combinações lineares interpretáveis das originais, como as da PCA, e
nem sequer são determinadas de forma única (rodar de novo com outra semente dá
outro mapa). Alimentar um classificador com elas é usar como covariável o
resultado de um algoritmo de visualização.

A PCA, ao contrário, é uma transformação linear fixa: ela aprende uma matriz no
treino e a aplica a qualquer ponto novo. É por isso que ela cabe num
\texttt{Pipeline} e o t-SNE não.
\end{solucao}

%------------------------------------------------------------------------------
\begin{exercicio}[PCA e SVD são a mesma coisa]
Seja $\bm X$ a matriz $n\times d$ dos dados \textbf{centrados}, com decomposição
em valores singulares $\bm X = \bm U\bm D\bm V^\top$, em que $\bm U$ é
$n\times r$ com colunas ortonormais, $\bm D$ é diagonal $r\times r$ com
$d_1\ge d_2\ge\dots\ge d_r>0$, e $\bm V$ é $d\times r$ com colunas ortonormais.
\begin{enumerate}[label=(\alph*),leftmargin=1.1cm]
  \item Mostre que as colunas de $\bm V$ são autovetores de
        $\bm S=\tfrac1n\bm X^\top\bm X$ e que o autovalor associado a $\bm v_j$
        é $\lambda_j = d_j^2/n$.
  \item Conclua que as componentes principais podem ser calculadas por SVD, sem
        formar $\bm X^\top\bm X$. Por que isso importa numericamente?
  \item As coordenadas dos dados nas $k$ primeiras componentes são as $k$
        primeiras colunas de $\bm U\bm D$. Verifique isso a partir de
        $\bm Z = \bm X\bm V$.
  \item O teorema de Eckart--Young diz que
        $\bm X_k = \bm U_k\bm D_k\bm V_k^\top$ é a melhor aproximação de posto
        $k$ de $\bm X$, no sentido de minimizar $\norm{\bm X - \bm A}_F$ entre
        todas as $\bm A$ de posto $\le k$. Que interpretação isso dá à PCA?
\end{enumerate}
\end{exercicio}

\begin{solucao}
\textbf{(a)} Usando $\bm U^\top\bm U = \bm I$,
\[
  \bm X^\top\bm X = (\bm U\bm D\bm V^\top)^\top(\bm U\bm D\bm V^\top)
   = \bm V\bm D\,\underbrace{\bm U^\top\bm U}_{\bm I}\,\bm D\bm V^\top
   = \bm V\bm D^2\bm V^\top .
\]
Multiplicando à direita por $\bm v_j$ (a $j$-ésima coluna de $\bm V$) e usando
$\bm V^\top\bm v_j = \bm e_j$:
\[
  \bm X^\top\bm X\,\bm v_j = \bm V\bm D^2\bm e_j = d_j^2\,\bm v_j .
\]
Dividindo por $n$, $\bm S\bm v_j = (d_j^2/n)\,\bm v_j$: os $\bm v_j$ são
autovetores de $\bm S$ com autovalores $\lambda_j = d_j^2/n$. Como
$d_1\ge\dots\ge d_r$, eles já vêm em ordem decrescente de autovalor.

\smallskip
\textbf{(b)} Basta calcular a SVD de $\bm X$ e ler $\bm V$ e $\bm D$; nunca é
preciso formar $\bm X^\top\bm X$.

Isso importa porque formar $\bm X^\top\bm X$ \textbf{eleva ao quadrado o número
de condição} da matriz. Se $\bm X$ tem valores singulares indo de $d_1$ a $d_r$,
o número de condição de $\bm X$ é $d_1/d_r$ e o de $\bm X^\top\bm X$ é
$(d_1/d_r)^2$. Em precisão dupla, com $\kappa(\bm X)\approx 10^8$ --- nada
absurdo em dados reais correlacionados --- o produto tem $\kappa\approx 10^{16}$
e as componentes menores viram ruído numérico.

É o mesmo motivo pelo qual bons solucionadores de mínimos quadrados usam QR ou
SVD em vez de resolver as equações normais diretamente (Aula~02).

\smallskip
\textbf{(c)} As coordenadas nas componentes principais são as projeções sobre os
$\bm v_j$, isto é $\bm Z = \bm X\bm V$. Substituindo a SVD e usando
$\bm V^\top\bm V=\bm I$:
\[
  \bm Z = \bm X\bm V = \bm U\bm D\bm V^\top\bm V = \bm U\bm D .
\]
As $k$ primeiras colunas de $\bm Z$ são, portanto, as $k$ primeiras de
$\bm U\bm D$ --- ou seja, $\bm u_j d_j$ para $j=1,\dots,k$. (Confere com o item
(a): a variância da coluna $j$ é $\tfrac1n\norm{\bm u_j d_j}^2 = d_j^2/n =
\lambda_j$, já que $\norm{\bm u_j}=1$.)

\smallskip
\textbf{(d)} Dá à PCA uma caracterização que não menciona variância nenhuma:
\textbf{a projeção sobre as $k$ primeiras componentes principais é a
aproximação de posto $k$ que menos distorce a matriz de dados}, medida pela soma
dos quadrados de todos os erros de reconstrução.

São duas leituras equivalentes do mesmo objeto:
\begin{itemize}[leftmargin=1.4cm]
  \item \emph{maximizar} a variância retida nas $k$ direções escolhidas;
  \item \emph{minimizar} o erro quadrático de reconstrução ao voltar de $k$
        dimensões para $d$.
\end{itemize}
E elas são equivalentes pelo mesmo mecanismo do Exercício~4 da Lista Teórica~E1:
a soma de quadrados total é fixa, então maximizar a parte ``explicada'' é
minimizar a parte ``residual''. O $k$-médias parte a soma de quadrados entre
grupos e dentro de grupos; a PCA a parte entre direções retidas e direções
descartadas. É a mesma identidade, aplicada a dois tipos de estrutura.
\end{solucao}


\end{document}
